\chapter{Composite Bound States and Bethe--Salpeter Amplitudes}
\label{ch:volume-ii-composite-bound-states-bethe-salpeter}

\section{Composite Particles as Poles}

Let \(\phi\) be a real scalar field whose stable one-particle states have
mass \(m\). A stable two-particle composite state \(B\) of mass \(M_B\) is a
normalizable one-particle state in the spectrum with
\[
  0<M_B<2m .
\]
The inequality places the state below the two-\(\phi\) continuum threshold.
It is then represented on the first sheet of the appropriate four-point
function by a pole at
\[
  s=M_B^2,
  \qquad
  s=-(k_1+k_2)^2 .
\]
After LSZ reduction of the external stable \(\phi\) particles, the same pole
appears in the connected scattering amplitude.

\begin{center}
\begin{tikzpicture}[>=Stealth,scale=0.95]
  \begin{scope}
    \draw[->,line width=0.75pt] (-0.35,0)--(6.7,0) node[right] {\(s\)};
    \draw[->,line width=0.75pt] (0,-1.35)--(0,1.9)
      node[above] {\(\operatorname{Im}s\)};
    \fill[blue!75!black] (2.15,0) circle (2.5pt);
    \node[blue!75!black,below] at (2.15,-0.18) {\(M_B^2\)};
    \draw[red!75!black,line width=1.8pt] (4.3,0)--(6.45,0);
    \node[red!75!black,below] at (4.3,-0.18) {\(4m^2\)};
    \node[blue!75!black,align=center] at (2.0,1.0)
      {first-sheet\\composite pole};
    \draw[blue!75!black,->,line width=0.8pt] (2.0,0.68)--(2.12,0.12);
    \node[red!75!black,align=center] at (5.45,0.72)
      {two-particle\\cut};
  \end{scope}
  \begin{scope}[xshift=9.1cm]
    \node[draw,circle,fill=gray!15,minimum size=0.92cm,inner sep=0pt] (b) at (0,0) {};
    \draw[line width=0.85pt] (-1.35,0.82)--(b.145);
    \draw[line width=0.85pt] (-1.35,-0.82)--(b.215);
    \draw[line width=0.85pt] (b.35)--(1.35,0.82);
    \draw[line width=0.85pt] (b.325)--(1.35,-0.82);
    \node at (0,0) {\(\mathcal M\)};
    \node[left] at (-1.35,0.82) {\(k_1\)};
    \node[left] at (-1.35,-0.82) {\(k_2\)};
    \node[right] at (1.35,0.82) {\(k_3\)};
    \node[right] at (1.35,-0.82) {\(k_4\)};
  \end{scope}
\end{tikzpicture}
\end{center}

This statement is independent of whether the Lagrangian contains a field
named \(B\). What matters is the spectral state and the residue with which
operators in the two-particle channel create it.

\section{A Bubble-Chain Example}

Consider the scalar theory
\[
  \mathcal L
  =
  -\frac12(\partial_\mu\phi)^2
  -\frac12m^2\phi^2
  -\frac{g}{4!}\phi^4 .
\]
The four-point vertex is \(-\ii g\). In the \(s\)-channel bubble-chain
approximation to \(2\to2\) scattering, the one-loop bubble is
\[
  \frac{(-\ii g)^2}{2}
  \int\frac{\dd^D\ell}{(2\pi)^D}\,
  \frac{-\ii}{\ell^2+m^2-\ii0}\,
  \frac{-\ii}{(k_1+k_2-\ell)^2+m^2-\ii0}.
\]
The factor \(1/2\) is the symmetry factor of the two identical internal
lines. Feynman parametrization and Wick rotation give
\[
  \frac{\ii g^2}{2}
  \int_0^1\dd x
  \int\frac{\dd^D\ell_E}{(2\pi)^D}\,
  \frac{1}{\left[\ell_E^2-sx(1-x)+m^2-\ii0\right]^2}.
\]

In \(D=2\), this becomes
\[
  \frac{\ii g^2}{8\pi} f(s),
  \qquad
  f(s)=
  \int_0^1\dd x\,
  \frac{1}{m^2-sx(1-x)-\ii0}.
\]
The geometric bubble chain is then
\[
  \ii\mathcal M_{\mathrm{chain}}(s)
  =
  -\ii g
  \sum_{L=0}^{\infty}
  \left(-\frac{g}{8\pi}f(s)\right)^L
  =
  \frac{-\ii g}{1+\dfrac{g}{8\pi}f(s)} .
\]
External wavefunction factors multiply the expression and do not move the
zeros of the denominator.

\begin{center}
\begin{tikzpicture}[>=Stealth,scale=0.82,
  leg/.style={line width=0.8pt},
  dot/.style={fill=gray!25,draw,circle,inner sep=1.8pt}]
  \begin{scope}
    \node at (-0.9,0) {\(\ii\mathcal M_{\mathrm{chain}}=\)};
    \node[dot] (c) at (0.35,0) {};
    \draw[leg] (-0.35,0.55)--(c)--(1.05,0.55);
    \draw[leg] (-0.35,-0.55)--(c)--(1.05,-0.55);
    \node at (1.55,0) {\(+\)};
    \begin{scope}[xshift=2.6cm]
      \draw[leg] (-0.75,0.55)--(-0.2,0.16);
      \draw[leg] (-0.75,-0.55)--(-0.2,-0.16);
      \draw[leg] (0.2,0.16)--(0.75,0.55);
      \draw[leg] (0.2,-0.16)--(0.75,-0.55);
      \draw[leg] (0,0) circle (0.34);
      \node at (1.12,0) {\(+\)};
    \end{scope}
    \begin{scope}[xshift=4.9cm]
      \draw[leg] (-0.95,0.55)--(-0.48,0.16);
      \draw[leg] (-0.95,-0.55)--(-0.48,-0.16);
      \draw[leg] (0.48,0.16)--(0.95,0.55);
      \draw[leg] (0.48,-0.16)--(0.95,-0.55);
      \draw[leg] (-0.26,0) circle (0.30);
      \draw[leg] (0.26,0) circle (0.30);
      \node at (1.38,0) {\(+\cdots\)};
    \end{scope}
    \node[align=center,blue!70!black,font=\small] at (3.25,-1.25)
      {pole condition:\\\(1+g f(s)/(8\pi)=0\)};
  \end{scope}
\end{tikzpicture}
\end{center}

Near the two-particle threshold from below, write
\[
  s=4m^2-\delta,
  \qquad
  \delta\to0^+ .
\]
The integral is dominated by \(x=1/2\). Setting \(x=1/2+y\), one obtains
\[
  m^2-sx(1-x)
  =
  \frac{\delta}{4}+4m^2y^2+O(\delta y^2),
\]
and hence
\[
  f(4m^2-\delta)
  \simeq
  \frac{\pi}{m\sqrt{\delta}} .
\]
The denominator can vanish for attractive coupling \(g<0\). At weak
coupling,
\[
  \delta
  \simeq
  \left(\frac{g}{8m}\right)^2,
  \qquad
  M_B
  =
  \sqrt{4m^2-\delta}
  \simeq
  2m-\frac{g^2}{256m^3}.
\]
Thus the chain produces a stable pole slightly below the two-particle
threshold.

The calculation uses the local expansion around a perturbative vacuum. A
complete scalar model with an attractive quartic term must specify the
large-field behavior of the potential. For example,
\[
  V(\phi)=\frac{m^2}{\beta^2}\left(1-\cos\beta\phi\right)
  =
  \frac12m^2\phi^2
  -\frac{m^2\beta^2}{4!}\phi^4
  +O(\phi^6)
\]
has an attractive quartic term in its local expansion while retaining a
well-defined periodic potential.

\section{Bethe--Salpeter Amplitude}

Let \(\ket{B(P)}\) be a stable composite one-particle state of mass \(M_B\),
with
\[
  P^2=-M_B^2,\qquad P^0>0 .
\]
The Bethe--Salpeter amplitude associated with the scalar field \(\phi\) is
the time-ordered matrix element
\[
  \Phi_B(x_1,x_2;P)
  =
  \langle\Omega|T\phi(x_1)\phi(x_2)|B(P)\rangle .
\]
Translation covariance gives
\[
  \Phi_B(x_1+a,x_2+a;P)
  =
  \ee^{-\ii P\cdot a}\Phi_B(x_1,x_2;P).
\]
Therefore its Fourier transform has support on total momentum \(P\):
\[
  \Phi_B(x_1,x_2;P)
  =
  \int\frac{\dd^D k_1}{(2\pi)^D}
       \frac{\dd^D k_2}{(2\pi)^D}\,
  \ee^{\ii k_1\cdot x_1+\ii k_2\cdot x_2}
  (2\pi)^D\delta^D(k_1+k_2-P)
  \widetilde\Phi_B(k_1,k_2;P).
\]
The normalization of \(\widetilde\Phi_B\) depends on the normalization of the
one-particle state; the pole residue fixes it.

Now consider the momentum-space time-ordered four-point function
\[
\begin{aligned}
  \widetilde G_4(k_1,k_2,k_3,k_4)
  &=
  \int\prod_{j=1}^4\dd^D x_j\,
  \exp\!\left[-\ii\sum_{j=1}^4 k_j\cdot x_j\right]
\\
  &\qquad\qquad\times
  \langle\Omega|T\phi(x_1)\phi(x_2)\phi(x_3)\phi(x_4)|\Omega\rangle .
\end{aligned}
\]
All four momenta are taken incoming, so translation invariance supplies
\((2\pi)^D\delta^D(k_1+k_2+k_3+k_4)\). Near the bound-state pole in the
\((12)\)-channel,
\[
  P=k_1+k_2,\qquad P^2\to -M_B^2,
\]
insertion of a complete set of states gives the factorization
\[
\begin{aligned}
  \widetilde G_4(k_1,k_2,k_3,k_4)
  &\sim
  (2\pi)^D\delta^D(k_1+k_2+k_3+k_4)
\\
  &\quad\times
  \frac{-\ii\,
    \widetilde\Phi_B(k_1,k_2;P)\,
    \widetilde\Phi_B^*(-k_3,-k_4;P)}
       {P^2+M_B^2-\ii0}
  +\text{terms regular at }P^2=-M_B^2 .
\end{aligned}
\]
This is the four-point analogue of the spectral pole in a two-point
function. The residue is the product of the amplitudes for two fields to
create and annihilate the stable composite state.

\begin{center}
\begin{tikzpicture}[>=Stealth,scale=0.9,
  leg/.style={line width=0.85pt},
  blob/.style={draw,fill=gray!15,circle,minimum size=0.92cm,inner sep=0pt},
  bs/.style={draw,fill=green!10,ellipse,minimum width=1.15cm,minimum height=1.0cm,inner sep=0pt}]
  \begin{scope}
    \node[blob] (g) at (0,0) {\(\widetilde G_4\)};
    \draw[leg] (-1.55,0.82)--(g.145);
    \draw[leg] (-1.55,-0.82)--(g.215);
    \draw[leg] (g.35)--(1.55,0.82);
    \draw[leg] (g.325)--(1.55,-0.82);
    \node[left] at (-1.55,0.82) {\(k_1\)};
    \node[left] at (-1.55,-0.82) {\(k_2\)};
    \node[right] at (1.55,0.82) {\(-k_3\)};
    \node[right] at (1.55,-0.82) {\(-k_4\)};
  \end{scope}
  \node at (3.0,0) {\(\sim\)};
  \begin{scope}[xshift=6.0cm]
    \node[bs] (l) at (-1.35,0) {\(\widetilde\Phi_B\)};
    \node[bs] (r) at (1.35,0) {\(\widetilde\Phi_B^*\)};
    \draw[leg] (-3.0,0.82)--(l.145);
    \draw[leg] (-3.0,-0.82)--(l.215);
    \draw[leg] (r.35)--(3.0,0.82);
    \draw[leg] (r.325)--(3.0,-0.82);
    \draw[blue!70!black,line width=1.5pt] (l.east)--(r.west);
    \node[blue!70!black,above] at (0,0.18)
      {\(\displaystyle \frac{-\ii}{P^2+M_B^2-\ii0}\)};
    \node[blue!70!black,below] at (0,-0.28) {\(B\)};
  \end{scope}
\end{tikzpicture}
\end{center}

\section{The Homogeneous Equation}

Let \(p\) and \(q\) denote relative momenta at fixed total momentum \(P\).
For equal constituent masses one may write the two internal momenta as
\[
  \frac{P}{2}+q,\qquad \frac{P}{2}-q .
\]
Let \(D(k)\) be the exact time-ordered propagator of the constituent field,
and define the two-particle propagator
\[
  \mathcal G_P(q)
  =
  D\!\left(\frac{P}{2}+q\right)
  D\!\left(\frac{P}{2}-q\right).
\]
Let \(\mathcal K_P(p,q)\) be the amputated two-particle irreducible kernel in
the chosen channel: by definition, it cannot be separated into two pieces by
cutting the two constituent lines carrying total momentum \(P\).

The amputated four-point scattering block \(\mathcal T_P(p,p')\) obeys the
operator equation
\[
  \mathcal T_P
  =
  \mathcal K_P
  +
  \mathcal K_P\mathcal G_P\mathcal T_P ,
\]
or explicitly
\[
\begin{aligned}
  \mathcal T_P(p,p')
  &=
  \mathcal K_P(p,p')
\\
  &\quad+
  \int\frac{\dd^D q}{(2\pi)^D}\,
  \mathcal K_P(p,q)\mathcal G_P(q)\mathcal T_P(q,p') .
\end{aligned}
\]
A bound-state pole at \(P^2=-M_B^2\) has the form
\[
  \mathcal T_P(p,p')
  \sim
  \frac{\Psi_B(p;P)\,\overline{\Psi}_B(p';P)}
       {P^2+M_B^2-\ii0}
  +\text{regular terms}.
\]
Substituting this pole form into the integral equation and equating residues
gives the homogeneous Bethe--Salpeter equation
\[
  \Psi_B(p;P)
  =
  \int\frac{\dd^D q}{(2\pi)^D}\,
  \mathcal K_P(p,q)\mathcal G_P(q)\Psi_B(q;P),
  \qquad
  P^2=-M_B^2 .
\]
The unamputated Bethe--Salpeter amplitude is related to the amputated vertex by
\[
  \widetilde\Phi_B(p;P)
  =
  \mathcal G_P(p)\Psi_B(p;P),
\]
up to the same convention-dependent one-particle normalization used in the
four-point residue.

\begin{center}
\begin{tikzpicture}[>=Stealth,scale=0.88,
  leg/.style={line width=0.82pt},
  kern/.style={draw,fill=blue!8,rounded corners=2pt,minimum width=0.7cm,minimum height=1.25cm,inner sep=0pt},
  amp/.style={draw,fill=green!10,ellipse,minimum width=1.18cm,minimum height=1.0cm,inner sep=0pt}]
  \begin{scope}[xshift=-4.1cm]
    \node[kern] (k) at (0,0) {\(\mathcal K\)};
    \draw[leg] (-1.1,0.42)--(k.west|-0,0.42);
    \draw[leg] (-1.1,-0.42)--(k.west|-0,-0.42);
    \draw[leg] (k.east|-0,0.42)--(1.1,0.42);
    \draw[leg] (k.east|-0,-0.42)--(1.1,-0.42);
    \node[align=center,font=\small] at (0,-1.08)
      {two-particle\\irreducible};
  \end{scope}
  \node[align=center] at (-0.45,0.22)
    {\(\mathcal T_P=\mathcal K_P+\mathcal K_P\mathcal G_P\mathcal T_P\)};
  \node[align=center,font=\small] at (-0.45,-0.62)
    {inhomogeneous\\scattering equation};
  \begin{scope}[xshift=3.35cm]
    \node[amp] (psi) at (-1.25,0) {\(\Psi_B\)};
    \node at (-0.25,0) {\(=\)};
    \node[kern] (k2) at (0.72,0) {\(\mathcal K\)};
    \draw[leg] (1.08,0.42)--(1.95,0.42);
    \draw[leg] (1.08,-0.42)--(1.95,-0.42);
    \node[amp] (psi2) at (2.75,0) {\(\Psi_B\)};
    \draw[leg] (1.95,0.42)--(psi2.west|-0,0.25);
    \draw[leg] (1.95,-0.42)--(psi2.west|-0,-0.25);
    \node[font=\small,blue!70!black] at (1.68,0.78) {\(\mathcal G_P\)};
    \node[align=center,font=\small] at (1.0,-1.08)
      {zero mode at\\\(P^2=-M_B^2\)};
  \end{scope}
\end{tikzpicture}
\end{center}

\section{Ladder Limit}

The kernel \(\mathcal K_P\) is exact when it contains all two-particle
irreducible contributions in the chosen channel. A ladder approximation keeps
only a selected low-order kernel and iterates it through \(\mathcal G_P\). For
example, a single exchanged particle produces a kernel of the schematic form
\[
  \mathcal K_P(p,q)
  \approx
  \frac{g_{\mathrm{ex}}^2}
       {(p-q)^2+\mu^2-\ii0},
\]
with \(\mu\) the exchanged mass and with numerator factors determined by the
spins and couplings of the fields.

In a weakly bound massive system, let
\[
  P^0=2m+E_{\mathrm{nr}},
  \qquad |E_{\mathrm{nr}}|\ll m,
  \qquad |\vec p\,|\ll m .
\]
If the kernel is instantaneous to leading order, the relative-energy integral
through the two constituent propagators reduces the Bethe--Salpeter equation
to a Schrödinger eigenvalue problem for the equal-time wavefunction. This
nonrelativistic limit is a derived approximation to the four-point pole
condition. The relativistic equation retains the full dependence on relative
time, energy flow, and field-theoretic kernel structure.
